\documentclass[11pt]{article}
\usepackage[a4paper,margin=27mm]{geometry}
\usepackage{amsmath,amssymb,amsthm,xcolor,booktabs,array}
\newcommand{\PROVED}{\textcolor{green!40!black}{\textbf{PROVED}}}
\newcommand{\OPEN}{\textcolor{red!70!black}{\textbf{OPEN}}}
\newcommand{\AUDIT}{\textcolor{orange!70!black}{\textbf{AUDIT}}}
\newtheorem{theorem}{Theorem}
\newtheorem{corollary}{Corollary}
\title{Gamma state shorting and rational Euler leakage (v121)}
\date{13 August 2026}

\begin{document}
\maketitle

\section{A causal state factor for one Gamma cell}

For $a>0$ and $f\in C_c^\infty(\mathbb R)$ let $z_a$ be the causal state
\begin{equation}
 z_a'=-az_a+\sqrt{2a}\,f,\qquad z_a(-\infty)=0,              \tag{1}
\end{equation}
and define the all-pass output
\begin{equation}
 g_a=\sqrt{2a}\,z_a-f.                                      \tag{2}
\end{equation}

\begin{theorem}[Lossless Gamma state and difference energy --- \PROVED]
The state satisfies
\begin{align}
 \frac d{dt}|z_a|^2&=|f|^2-|g_a|^2,                         \tag{3}\\
 f-g_a&=\sqrt{\frac2a}\,z_a'.                               \tag{4}
\end{align}
Consequently
\begin{equation}
 \boxed{
 \frac1{2a}\|f-g_a\|_2^2
 =\frac1{a^2}\int_{\mathbb R}|z_a'(t)|^2\,dt
 =\int_{\mathbb R}
   \frac{2\tau^2}{a(a^2+\tau^2)}
   |\widehat f(\tau)|^2\frac{d\tau}{2\pi}.}                 \tag{5}
\end{equation}
\end{theorem}

\begin{proof}
Equations (1)--(2) give
\[
 |g_a|^2=2a|z_a|^2+|f|^2-2\sqrt{2a}\Re(\overline z_af),
\]
which is (3).  Solving (1) for $f$ and substituting in (2) gives (4).
In Fourier coordinates the transfer function of $f-g_a$ has squared
modulus $4\tau^2/(a^2+\tau^2)$.  This proves (5).
\end{proof}

For $a_j=2j+\frac12$, summing (5) gives the complete positive Gamma form:
\begin{equation}
 \langle G_\Gamma f,f\rangle
 =\sum_{j\ge0}\frac1{a_j^2}\int|z_{a_j}'(t)|^2\,dt.          \tag{6}
\end{equation}
This is the state-space version of the Gram (6)--(7) in v120.

\section{Exact shorting across a gap}

Let $J=(r,r+L)$ be a gap on which the physical source is free.  For one
Gamma state prescribe $z_a(r)=\zeta_0$ and $z_a(r+L)=\zeta_1$.

\begin{theorem}[Gamma gap resistance --- \PROVED]
Among all absolutely continuous state paths with those endpoint values,
\begin{equation}
 \inf\frac1{a^2}\int_r^{r+L}|z_a'(t)|^2\,dt
 =\frac{|\zeta_1-\zeta_0|^2}{a^2L}.                         \tag{7}
\end{equation}
The minimizer is the affine state path.  Its corresponding physical source,
obtained from (1), is
\begin{equation}
 f(t)=\frac{z_a'(t)+az_a(t)}{\sqrt{2a}}.                    \tag{8}
\end{equation}
\end{theorem}

\begin{proof}
Cauchy--Schwarz gives
\[
 |\zeta_1-\zeta_0|^2
 =\left|\int_r^{r+L}z_a'(t)\,dt\right|^2
 \le L\int_r^{r+L}|z_a'|^2.
\]
Equality holds exactly for constant $z_a'$, proving (7); (8) is (1).
\end{proof}

For the full Gamma place the same physical source $f$ drives every state
$z_{a_j}$.  If this common-control constraint is relaxed and each channel
is allowed its own gap source, (7) gives the rigorous lower bound
\begin{equation}
 \operatorname{Short}_{J}(G_\Gamma)
 \ge
 \sum_{j\ge0}\frac{|\zeta_{j,1}-\zeta_{j,0}|^2}{a_j^2L}.     \tag{9}
\end{equation}
Indeed, relaxation enlarges the class over which the positive energy is
minimized.  Thus the right side is an independently positive boundary
network; it is not obtained by taking a square root of the desired Schur
complement.

\section{Integer incidence and rational leakage}

On the right-oriented threshold fiber label the interval
$T_N-\log m+(0,\delta_N)$ by the integer $m\le N$.  Translation by
$+\log n$ on the source corresponds to multiplication of labels,
\begin{equation}
                         m\longmapsto mn,                     \tag{10}
\end{equation}
and remains in the integer fiber exactly when $mn\le N$.  Translation in
the opposite direction corresponds geometrically to
\begin{equation}
                         m\longmapsto m/n.                    \tag{11}
\end{equation}
It lies in the integer fiber exactly when $n\mid m$.

\begin{theorem}[Exact decomposition of the Euler boundary action ---
\PROVED]
On the arithmetic threshold fiber, every finite Euler translation splits
orthogonally into
\begin{equation}
 S_{\log n}=S_{\log n}^{\rm div}\oplus S_{\log n}^{\rm leak},\tag{12}
\end{equation}
where the first part is the finite divisibility shift (10), together with
its adjoint on divisible labels, and the second part has support on the
nonintegral rational labels $m/n$ with $n\nmid m$.  Distinct rational
labels whose logarithmic cells do not overlap give orthogonal leakage
channels.  Every possible overlap is determined by the explicit condition
\begin{equation}
 \left|\log\frac{m/n}{m'/n'}\right|<\delta_N.                 \tag{13}
\end{equation}
\end{theorem}

\begin{proof}
Translation of the interval positions gives (10)--(11).  A noninteger
label $m/n$ cannot overlap an integer label $k$ in positive measure.
Indeed, if $m>kn$, then the weakest possible separation occurs at the
largest numerator and
\[
 \frac{m}{kn}\ge\frac{N}{N-1}>\frac{N+1}{N}=e^{2\delta_N}.
\]
If $m<kn$, then $kn/m\ge(N+1)/N=e^{2\delta_N}$, with equality
giving only an endpoint contact.  Thus the integer and noninteger support
unions are orthogonal.  Two intervals of length $\delta_N$ overlap
precisely when the distance between their logarithmic left endpoints is
less than $\delta_N$, which is (13).  The integral part is exactly the
truncated divisibility representation used in G39.
\end{proof}

\paragraph{Scope audit.}
The word Grams of v109--v110 and the cross transport of v119 control
$S^{\rm div}$.  Equations (6)--(9) construct a positive Gamma boundary
network for the gaps.  The only Euler component not yet assigned to that
network is the clustered rational leakage in (13).  Calling all leakage
channels orthogonal would be false because distinct fractions can fall
within one cell width.

The sharp remaining cross-place theorem can therefore be stated without a
global inverse:
\begin{equation}
 \boxed{\text{the Gram of the clustered rational leakage is dominated,
 gap by gap, by the Gamma resistance (9) plus the Tate terminal of v117}.}
                                                                    \tag{14}
\end{equation}
If (14) holds with unit constant, elimination of every gap is passive;
cascading the cells proves the shorted domination v120 (10), hence (8c).

\paragraph{Cascade audit after v122 --- \AUDIT.}
The local boundary-state estimate v122 (12) has a strict numerical margin,
but it may not be applied by resetting the incoming Gamma state to zero on
each rational cell.  In a connected leakage cluster the state exiting one
cell enters the next.  The signed endpoint term in v122 (8) telescopes,
whereas the transported nearly constant state survives.  Therefore the
local estimate proves the one-cell gate but not (14) for an arbitrarily
long cluster.  Closing (14) requires a multiscale estimate for the complete
tower $a_j=2j+\frac12$, or an equivalent global discrete Poincar\'e
inequality including the two outer Tate states.  No reset assumption is
used in the ledger.

\begin{center}
\begin{tabular}{>{\raggedright\arraybackslash}p{.55\linewidth}
                >{\centering\arraybackslash}p{.15\linewidth}
                >{\raggedright\arraybackslash}p{.20\linewidth}}
\toprule
Statement & Status & Location \\
\midrule
Causal lossless Gamma factor & \PROVED & (1)--(6) \\
Exact Gamma gap resistance & \PROVED & (7)--(9) \\
Divisibility/rational-leakage split & \PROVED & (10)--(13) \\
Unit domination of clustered leakage by Gamma--Tate resistance & \OPEN & (14) \\
Contractivity (8c), condition $G$, row (d) & \OPEN & after (14) \\
\bottomrule
\end{tabular}
\end{center}

\end{document}
